\documentclass[11pt,a4paper]{article}
\usepackage[utf8]{inputenc}
\usepackage[T1]{fontenc}
\usepackage{amsmath,amsfonts,amssymb}
\usepackage{graphicx}
\usepackage{hyperref}
\usepackage{listings}
\usepackage{color}
\usepackage{booktabs}
\usepackage{float}
\usepackage{geometry}
\geometry{margin=1in}
\definecolor{zenblue}{RGB}{41,121,255}
\hypersetup{colorlinks=true,linkcolor=zenblue,urlcolor=zenblue,citecolor=zenblue}

\title{\textbf{Zen3-Guard: A Packaging of IBM Granite Guardian 8B\\
for the Zen Model Family}\\
\large Technical Whitepaper v2025.06}
\author{Zach Kelling \\ Zen LM Research Team\\
\texttt{research@zenlm.org}\\
\href{https://papers.zenlm.org}{papers.zenlm.org}}
\date{June 2025}

\begin{document}
\maketitle

\begin{abstract}
Zen3-Guard is a repackaging of IBM's \textbf{Granite Guardian 8B} safety model, deployed
as the content-safety layer for the Zen model family. It is \emph{not} a from-scratch
model: it is the openly released, Apache-2.0 licensed \texttt{ibm-granite/granite-guardian}
8B guardrail produced by IBM Research, which is itself a supervised fine-tune of the dense,
decoder-only IBM Granite 3.x 8B Instruct base model (\texttt{GraniteForCausalLM}). Granite
Guardian classifies a user prompt and/or assistant response against a configurable risk
definition and emits a \texttt{Yes}/\texttt{No} verdict together with a probability of harm
derived from the next-token logits. The risk taxonomy, aligned with the IBM AI Risk Atlas,
spans harmful content (harm, social bias, profanity, sexual content, unethical behavior,
violence, jailbreaking) as well as retrieval-augmented-generation faithfulness risks
(context relevance, groundedness, answer relevance). IBM reports an aggregate AUC of
\textbf{0.871} on harmful-content benchmarks and \textbf{0.854} on RAG-hallucination
benchmarks for Granite Guardian 3.0 8B~\cite{graniteguardian}. This document describes how
the upstream model is integrated as a drop-in safety filter for Zen deployments; all model
weights, training, and reported metrics are IBM's and are cited as such. We add no
benchmark claims of our own.
\end{abstract}

\tableofcontents
\newpage

%% -----------------------------------------------------------------------
\section{Introduction}
\label{sec:intro}

Content safety classification is a critical infrastructure component for any deployed
language model service. Effective moderation must satisfy competing constraints:
high recall on harmful content (to prevent policy violations), high precision (to
avoid false positives that degrade user experience), coverage across 100+ languages
(to serve global user populations), and sub-10ms inference latency (to insert into
API request paths without perceptible overhead).

Existing safety approaches trade off along several dimensions:
\begin{itemize}
\item Keyword filters: high false positive rate, trivially bypassed by paraphrasing
\item Narrow binary classifiers: insufficient granularity for nuanced policy enforcement
\item Large safety-tuned generation models: accurate but heavier to serve
\item Embedding-similarity approaches: brittle to adversarial rephrasing
\end{itemize}

Rather than train a safety classifier from scratch, Zen3-Guard adopts IBM's openly
released Granite Guardian 8B~\cite{graniteguardian}, which directly targets these
trade-offs as a purpose-built guardrail LLM. Our contribution is integration and
packaging, not modeling. The upstream model provides:

\begin{enumerate}
\item \textbf{A risk-conditioned guardrail LLM} --- A dense Granite 3.x 8B Instruct model
      fine-tuned by IBM to classify content against a named risk definition supplied at
      inference time, emitting a \texttt{Yes}/\texttt{No} verdict plus a harm probability.
\item \textbf{An IBM AI Risk Atlas taxonomy} --- Risk dimensions covering harmful content
      and (in the RAG setting) faithfulness risks, rather than a single binary safe/unsafe
      label.
\item \textbf{Prompt- and response-side checking} --- The same model can be applied to the
      user input, the model output, or a RAG response against its retrieved context.
\item \textbf{Risk-definition prompting} --- Because the desired risk is described in the
      prompt template at inference time, the deployed model can be steered toward different
      checks without retraining.
\end{enumerate}

We do not modify the Granite Guardian weights; Zen3-Guard is the upstream model wired into
the Zen serving stack. The remainder of this document describes that integration and
summarizes IBM's published design and results with attribution.

%% -----------------------------------------------------------------------
\section{Architecture}
\label{sec:arch}

\subsection{Upstream Model: Granite Guardian 8B}

Zen3-Guard is the IBM Granite Guardian 8B model served unmodified. Per IBM's
documentation~\cite{graniteguardian}, the model is a \emph{dense, decoder-only}
transformer (\texttt{GraniteForCausalLM}; GQA, RoPE, SwiGLU MLP, RMSNorm). It is
\emph{not} a mixture-of-experts model and has no notion of ``expert clusters'' or
safety-specific routing. It is produced by supervised fine-tuning of the IBM Granite 3.x
8B Instruct base model on risk-annotated data. We make no architectural changes; earlier
revisions of this document incorrectly described a bespoke ``mixture-of-distilled-experts''
backbone, which does not correspond to the deployed weights.

\subsection{Risk Taxonomy and Output Schema}

Granite Guardian does not emit a fixed multi-label vector. Instead, the caller supplies a
\emph{risk definition} in the prompt template, and the model judges whether the supplied
content exhibits that risk. The risk dimensions documented by IBM~\cite{graniteguardian},
aligned with the IBM AI Risk Atlas, include:

\begin{itemize}
\item \textbf{Harmful content}: harm (general), social bias, profanity, sexual content,
      unethical behavior, violence, and jailbreaking.
\item \textbf{RAG faithfulness}: context relevance, groundedness, and answer relevance
      (i.e., detecting ungrounded or off-topic answers).
\end{itemize}

The IBM AI Risk Atlas mapping organizes harm into a deeper hierarchy (four high-level
categories and thirteen sub-categories), but the model's runtime interface is per-risk
rather than a flat 47-way classifier.

For the selected risk, the model produces:
\begin{itemize}
\item Binary verdict: a \texttt{Yes} / \texttt{No} token (\texttt{Yes} = risk present).
\item Probability of harm: derived from the log-probabilities of the \texttt{Yes}/\texttt{No}
      tokens at the final position.
\end{itemize}

Severity tiers and character-offset evidence spans are not native outputs of the upstream
model; where downstream Zen tooling needs them, they are derived heuristically outside the
model and are not part of Granite Guardian's reported behavior.

\subsection{Language Coverage}

Granite Guardian inherits the tokenizer and multilingual coverage of its Granite 3.x base.
IBM's published evaluations focus primarily on English safety benchmarks; we therefore do
not assert specific per-language safety scores. Claims in earlier revisions of a custom
150{,}528-token multilingual tokenizer and a cross-lingual transfer objective were
fabricated and have been removed.

\subsection{Risk-Definition Prompting}

Because the risk to be checked is described in the chat template rather than baked into a
classification head, a single deployed instance can perform different checks (e.g.
\texttt{harm}, \texttt{social\_bias}, \texttt{groundedness}) by changing the
\texttt{risk\_name} passed to the tokenizer's \texttt{apply\_chat\_template}. Schematically:

\begin{verbatim}
apply_chat_template(
    messages=[{"role": "user", "content": <user content>},
              {"role": "assistant", "content": <response>}],
    guardian_config={"risk_name": "harm"})
# -> model emits "Yes"/"No"; prob_of_risk read from token logits
\end{verbatim}

This is the upstream interface; Zen deployments expose it through a thin wrapper. It is a
prompt-time mechanism, not a learned ``policy enforcement expert cluster.''

%% -----------------------------------------------------------------------
\section{Training}
\label{sec:training}

\subsection{Provenance, Not Training}

Zen3-Guard performs no training of its own. The guardrail was trained by IBM Research;
we summarize their published methodology for completeness and attribute it
accordingly~\cite{graniteguardian}.

Per IBM, Granite Guardian is created by \emph{supervised fine-tuning} of the Granite 3.x
8B Instruct base model on a combination of human-annotated and synthetic risk data
spanning the harm and RAG-faithfulness dimensions of the IBM AI Risk Atlas. The objective
is standard next-token supervision over the \texttt{Yes}/\texttt{No} verdict given a
risk-conditioned prompt; the harm probability used downstream is read from the
\texttt{Yes}/\texttt{No} token log-probabilities at inference time.

The fabricated training pipeline in earlier revisions of this document --- a 12B-token
multilingual safety corpus, focal-loss ``safety expert'' specialization, a five-level
taxonomy hierarchy loss, and a four-phase adversarial regimen --- did not describe any real
training run and has been removed. For the authoritative dataset composition and training
recipe, see IBM's paper~\cite{graniteguardian}.

\subsection{Calibration}

We do not recalibrate the model. IBM reports its evaluation metrics (Section~\ref{sec:eval})
at a default decision threshold of $0.5$ on the harm probability; deployments may pick a
different threshold per risk to trade recall against false positives.

%% -----------------------------------------------------------------------
\section{Evaluation}
\label{sec:eval}

All numbers in this section are IBM's published results for Granite Guardian
3.0~8B~\cite{graniteguardian}, reproduced here with attribution. We ran no evaluations of
our own and make no independent benchmark claims.

\subsection{Aggregate Results Reported by IBM}

IBM reports two headline metrics for Granite Guardian 3.0~8B, aggregated over public
benchmark suites:

\begin{table}[H]
\centering
\caption{Granite Guardian 3.0 8B aggregate metrics, as reported by IBM~\cite{graniteguardian}.}
\begin{tabular}{lcc}
\toprule
\textbf{Capability} & \textbf{Metric} & \textbf{Value} \\
\midrule
Harmful-content detection      & Aggregate AUC & 0.871 \\
Harmful-content detection      & Aggregate F1 (threshold 0.5) & 0.758 \\
RAG hallucination / groundedness & Average AUC (TRUE suite) & 0.854 \\
\bottomrule
\end{tabular}
\label{tab:granite-agg}
\end{table}

The harmful-content figures are aggregated across public safety datasets used in IBM's
evaluation (for example ToxicChat, HarmBench, SafeRLHF, BeaverTails, OpenAI Moderation,
SimpleSafetyTests, and xstest), with comparisons against Llama~Guard and ShieldGemma
baselines; the RAG-hallucination figure is the average AUC on the TRUE
benchmark~\cite{graniteguardian}. We refer the reader to IBM's paper for per-dataset
breakdowns, which we do not reproduce here.

\subsection{A Note on the Numbers in Earlier Revisions}

Earlier revisions of this document reported figures such as ToxiGen 99.2\%, HatEval 97.8\%,
a multilingual macro-F1 of 0.968, sub-5ms A10G latency, and per-attack adversarial
robustness rates. None of these came from a measured evaluation of the deployed model; they
were fabricated and have been removed. Serving latency depends entirely on the operator's
hardware and stack and is not a property we can claim for the upstream weights.

%% -----------------------------------------------------------------------
\section{Related Work}
\label{sec:related}

\subsection{Content Moderation Systems}

Early content moderation relied on expert-curated keyword blocklists and regular
expression patterns. Machine learning approaches improved recall but suffered from
adversarial brittleness. Transformer-based classifiers~\cite{perspectiveapi} brought
contextual understanding but required separate models per language.

\subsection{Guardrail Models}

Constitutional AI~\cite{bai2022constitutional} demonstrated that safety behaviors can be
instilled in models through targeted fine-tuning. A subsequent line of work --- including
Llama~Guard, ShieldGemma, and IBM's Granite Guardian~\cite{graniteguardian} --- packages a
fine-tuned LLM specifically as an input/output guardrail that emits a risk verdict.
Zen3-Guard does not introduce a new model in this line; it adopts Granite Guardian directly.

\subsection{Multilingual Safety}

Cross-lingual safety classification has been studied through multilingual pre-training and
transfer~\cite{multilingual-safety}. We make no multilingual-safety claims for Zen3-Guard
beyond what IBM reports for Granite Guardian, whose published evaluations are primarily
English; the cross-lingual ``alignment mechanism'' described in earlier revisions did not
exist and has been removed.

%% -----------------------------------------------------------------------
\section{Deployment Recommendations}
\label{sec:deploy}

\subsection{Integration Patterns}

Zen3-Guard is designed for three primary deployment patterns:

\textbf{Synchronous API guard}: Insert as a pre-processing layer before the main
LLM, classifying the user prompt and rejecting or flagging it before generation.

\textbf{Parallel screening}: Run the guard concurrently with the main LLM on
the input, and suppress responses if it returns a positive risk verdict. On GPU clusters
with spare capacity this can overlap with generation.

\textbf{Output scanning}: Run the guard on model outputs (and, in RAG, against the
retrieved context for groundedness) before returning to the user. Recommended where
harmful or ungrounded content may emerge from a benign-looking input.

Because the guard is itself an 8B decoder LLM, its serving cost is comparable to running an
8B model; we make no specific latency guarantee, as it depends on the operator's hardware,
batching, and quantization.

\subsection{Threshold Configuration}

The harm probability is read from the \texttt{Yes}/\texttt{No} token log-probabilities.
IBM reports its metrics at a default threshold of $p = 0.5$~\cite{graniteguardian}.
Operators may raise the threshold per risk to favor precision (fewer false positives) or
lower it to favor recall (catching more borderline content); the appropriate operating
point is deployment-specific and should be chosen against the operator's own labeled data
rather than the illustrative recall/FPR figures asserted in earlier revisions, which were
not measured.

%% -----------------------------------------------------------------------
\section{Conclusion}
\label{sec:conclusion}

Zen3-Guard is the content-safety layer of the Zen model family, implemented as a direct
deployment of IBM's Apache-2.0 Granite Guardian 8B guardrail~\cite{graniteguardian}. It is
a dense, decoder-only Granite 3.x 8B Instruct model fine-tuned by IBM Research to classify
prompts and responses against a named risk and emit a \texttt{Yes}/\texttt{No} verdict with
a harm probability. IBM reports an aggregate AUC of 0.871 on harmful-content benchmarks and
0.854 on RAG-hallucination benchmarks for the 3.0 8B model. We contribute integration and
packaging into the Zen serving stack, not a new model, and we make no benchmark claims of
our own. Operators evaluating Zen3-Guard for production should validate it on their own
content distribution and choose per-risk thresholds accordingly.

\section*{Acknowledgments and Attribution}

Zen3-Guard is built entirely on IBM Granite Guardian (\texttt{ibm-granite/granite-guardian},
Apache-2.0), developed by IBM Research; all model weights, training data, and reported
metrics are IBM's. We thank the IBM Granite team for releasing the model openly.

\bibliographystyle{plain}
\begin{thebibliography}{9}

\bibitem{perspectiveapi}
E. Lees et al.,
``A New Generation of Perspective API: Efficient Multilingual Character-level Transformers,''
\textit{ACL}, 2022.

\bibitem{bai2022constitutional}
Y. Bai et al.,
``Constitutional AI: Harmlessness from AI Feedback,''
\textit{arXiv:2212.08073}, 2022.

\bibitem{multilingual-safety}
T. Hartvigsen et al.,
``ToxiGen: A Large-Scale Machine-Generated Dataset for Adversarial and Implicit Hate Speech Detection,''
\textit{ACL}, 2022.

\bibitem{graniteguardian}
I. Padhi et al. (IBM Research),
``Granite Guardian,''
\textit{arXiv:2412.07724}, 2024.
Model: \href{https://huggingface.co/ibm-granite/granite-guardian-3.0-8b}{ibm-granite/granite-guardian-3.0-8b} (Apache-2.0).

\end{thebibliography}

\end{document}
